\documentclass[11pt]{article}
\usepackage[a4paper,margin=25mm]{geometry}
\usepackage{amsmath}
\usepackage{graphicx}
\title{A Particle in a Box Obeys the Uncertainty Principle}
\author{BPhO Computational Physics Challenge 2026}
\date{}
\begin{document}
\maketitle

For an infinite one-dimensional box of width $a$, the normalised stationary
state is
\[
 \psi_n(x)=\sqrt{\frac{2}{a}}\sin\left(\frac{n\pi x}{a}\right),
 \qquad 0<x<a.
\]
The corresponding energy is $E_n=n^2\pi^2\hbar^2/(2ma^2)$.  The probability
density is symmetric about the centre of the box, therefore
\[
 \langle x\rangle=\frac{a}{2}.
\]
Direct integration gives
\[
 \langle x^2\rangle=a^2\left(\frac13-\frac{1}{2n^2\pi^2}\right),
\]
and consequently
\[
 \Delta x=a\sqrt{\frac1{12}-\frac{1}{2n^2\pi^2}}.
\]

The momentum expectation is zero because equal right- and left-moving parts
make up the standing wave. Applying $\hat p=-i\hbar\,d/dx$ twice yields
\[
 \langle p^2\rangle=\left(\frac{n\pi\hbar}{a}\right)^2,
 \qquad \Delta p=\frac{n\pi\hbar}{a}.
\]
Thus
\[
 \Delta x\Delta p=
 \frac{n\pi\hbar}{2}\sqrt{\frac13-\frac{1}{2n^2\pi^2}}
 =\frac{\hbar}{2}\sqrt{\frac{n^2\pi^2}{3}-1}.
\]
For $n=1$, this is $0.568\hbar$, already larger than the Heisenberg lower
bound $\hbar/2$.  It increases with $n$, so every stationary state obeys
$\Delta x\Delta p\geq\hbar/2$.  The accompanying program independently
checks the analytical result using numerical integration.

\begin{center}
\includegraphics[width=0.82\linewidth]{task7_uncertainty_principle.png}
\end{center}
\end{document}
